\secciongrande{Lenguajes de programación en entorno cliente dos entornos de ejecución}

\resumen{
  La programación web requiere diferenciar entre entornos cliente y servidor. El entorno cliente se ejecuta en el navegador del usuario, mientras que el entorno servidor opera en el servidor. La elección del lenguaje es crucial y depende de las ventajas e inconvenientes específicos, siendo el desarrollador quien toma la decisión final.
}

\definicion{Entorno cliente}{Entorno de ejecución donde el código es procesado por el dispositivo del usuario, típicamente un navegador web.}

\definicion{Entorno servidor}{Entorno de ejecución donde el código se procesa en el servidor remoto, gestionando datos y lógica de negocio.}

\subtituloA{Diferenciación de entornos en desarrollo web}

\lista{
  \textbf{Entorno cliente:} Ejecutado por el cliente en el navegador web, enfocado en la interacción con el usuario y la presentación. |
  \textbf{Entorno servidor:} Ejecutado en el servidor, responsable de la lógica de negocio, la gestión de bases de datos y la respuesta a peticiones.
}

\nota{Para el examen}{
Es vital comprender que cliente y servidor no son excluyentes; se complementan en el desarrollo de sitios web dinámicos. La elección del lenguaje debe basarse en sus pros y contras, no solo en la moda tecnológica.
}
